\documentclass[11pt]{article}
\usepackage[letterpaper,hmargin=1.02in,vmargin=0.93in]{geometry}
\usepackage[T1]{fontenc}
\usepackage{lmodern,amsmath,amssymb,amsthm,microtype}
\usepackage[hidelinks]{hyperref}
\newtheorem{theorem}{Theorem}
\newtheorem{lemma}[theorem]{Lemma}
\newtheorem{proposition}[theorem]{Proposition}
\newcommand{\F}{\mathbb F}
\newcommand{\NL}{\mathsf{NL}}
\newcommand{\US}{\mathsf{USPACE}}
\newcommand{\calC}{\mathcal C}
\newcommand{\calF}{\mathcal F}
\newcommand{\fr}{\operatorname{Fr}}
\newcommand{\id}{\operatorname{id}}
\newcommand{\pr}{\Pr}
\allowdisplaybreaks
\hypersetup{pdftitle={Unambiguous reachability in near-logarithmic space}}
\title{Unambiguous reachability in near-logarithmic space}
\author{}
\date{November 2026}
\begin{document}
\maketitle
\begin{abstract}
We give an unambiguous algorithm for directed reachability that uses
$O(\log n\,(\log\log n)^2)$ space. We construct a family of vertex weight assignments, indexed by short seeds,
that includes a min-isolating assignment for each layered graph.
The construction maintains a bound on the product of the numbers of optimal
boundary sequences for several endpoint pairs. This joint bound allows the
algorithm to double the number of binary levels handled at each step. We use
exhaustive seed search and make no polynomial-time claim.
\end{abstract}

\section{Statement and proof strategy}

\begin{theorem}\label{thm:main}
For every language in $\NL$ there is an unambiguous machine deciding the language
in $O(\log n\,(\log\log n)^2)$ space.
\end{theorem}

We give the proof for a simple layered DAG on $N$ vertices. Layers are numbered
$0,\ldots,d$ with $d\le N$, every edge advances one layer, and vertices have
distinct identifiers in $[N]$. The standard logspace reduction from $\NL$
gives graphs of this form with $N$ polynomial in the input length. Write
$L=\lceil\log_2 N\rceil$ and $\ell=\lceil\log_2(2L)\rceil$. All bounds below
are for sufficiently large $N$. If the graph contains no path of length two,
reachability can be decided directly in logarithmic space by testing equality
and adjacency. We henceforth assume at least one such path exists.

An assignment is \emph{min-isolating} if every reachable pair has a unique
minimum-weight path. We minimize a sequence of additive vertex weights
lexicographically. Optimal subpaths are optimal for their own endpoints,
by replacement. When comparing paths between fixed endpoints, we omit the
endpoints' contributions. The entire vector objective can be encoded by
one integer weight.

\paragraph{The issue with a single endpoint pair.}
Suppose that each binary midpoint has at most $K$ plausible choices.
Recording the path at $q-1$ boundaries gives as many as $K^{q-1}$
sequences for a fixed pair. Even after thinning to $K=\mathrm{poly}(L)$,
this estimate runs out when $q$ grows past a polynomial in $L$.
The proof below counts the choices in a collection of intervals jointly.

At a given scale, suppose a tuple of optimal boundary sequences has many
locations where two choices are possible. Choose $m$ such locations.
They involve disjoint sets of layers, so a random additive perturbation
must satisfy $m$ independent equations to preserve all the ties.
We need two bounds for the union argument. The first bounds how many
lists of locations and endpoints can occur. The second bounds how many
pairs of alternatives those locations admit, \emph{in total}. We build
the lists from coarse to fine using an initial midpoint bound. We then
choose weights from fine to coarse, proving the required joint bound
for every list. The lists are a counting device in the proof. The
machine only enumerates seeds for vertex weights.

The use of shifted hashes and the passage from min-isolation to unambiguous
computation follow van Melkebeek and Prakriya~\cite{vMP} and Reinhardt and
Allender~\cite{RA}. We use binary layer valuations in the form of the
dyadic construction analyzed by Hoza~\cite[Appendix A]{Hoza}.

\begin{lemma}[Small-bias labels]\label{lem:bias}
For a prime $p$, $n_0$ coordinates, and $0<\varepsilon<1$, there is an explicit
distribution on $h\in\F_p^{n_0}$ with every nontrivial additive character
of expectation of absolute value at most $\varepsilon$. A seed and the
space to evaluate a coordinate each use $O(\log n_0+\log p+\log(1/\varepsilon))$
bits. For any fixed system of $m$ linearly independent affine constraints
on $h$, the probability of satisfying all of them is at most
$p^{-m}+\varepsilon$.
\end{lemma}
\begin{proof}
Choose the smallest extension $E=\F_{p^u}$ with $|E|>n_0/\varepsilon$.
For uniform independent $x,y\in E$ and any nonzero
$\F_p$-linear functional $\pi:E\to\F_p$, put $h_j=\pi(yx^j)$, $1\le j\le n_0$.
A nontrivial linear combination has a nonzero coefficient polynomial
of degree at most $n_0$ in $x$. Away from its roots the linear combination
is uniform by the choice of $y$. This proves the bias bound.
Expanding the indicator of the affine system in characters gives a constant term
$p^{-m}$, and the remaining coefficients have total absolute value at most one.
For uniform evaluation, enumerate monic degree-$u$ polynomials and test
irreducibility by trial division. In the first resulting field representation
one may take $\pi$ to be the constant coefficient. The search and the field
arithmetic use $O(u\log p)$ space.
\end{proof}

\section{Thinning the feasible midpoints}

Write $L=\lceil\log_2 N\rceil$ and $\ell=\lceil\log_2(2L)\rceil$.
For two vertices $x,y$ in layers $a<b$ with $b-a\ge 2$, let their
\emph{pivot} be the internal layer $i$ maximizing $\nu_2(i)$.
The pivot is unique because between two positive integers of the same
$2$-adic valuation lies an integer of higher valuation. The feasible
midpoints for $(x,y)$ are the vertices in its pivot layer on an
$x$--$y$ path.

\begin{lemma}\label{lem:raw}
There is an explicitly enumerable family of nonnegative integer vertex
weights $V$ for $N$-vertex layered graphs with the following property.
For each graph, some $V$ in the family ensures that, for every reachable
pair $(x,y)$ having a pivot, at most $K=L^{O(1)}$ distinct pivot
vertices occur on minimum-$V$ paths from $x$ to $y$. The seed length,
the bit length of a path weight, and the workspace for evaluating a
vertex weight are all $O(L\ell^2)$.
\end{lemma}

\begin{proof}
We first construct a coloring for any fixed family of at most $N^2$
nonempty subsets of an $N$-element ground set. We want at most $K$
elements of the least occupied color in every set. At any stage retain
in each set the elements having its lexicographically least tuple of
colors so far. Suppose every retained set has size at most $2^a$.
Initially $a=L$. While $a>256\ell$, the next stage reduces this bound to
$2^{\lceil a/2\rceil}$.

Choose a fresh $\varepsilon$-biased string with $Na$ binary coordinates
using Lemma~\ref{lem:bias} over $\mathbb F_2$. Each vertex receives its
own block of $a$ bits. Its color for this stage is $a$ minus the number
of leading zero bits in its block. In particular, for every $d\le a$,
the event that the first $d$ bits are zero is a union of the lowest
colors. Take
\[
 t=2\lceil 64L/a\rceil,\qquad
 \varepsilon=2^{-10L-2t(a+1)}.
\]
This choice uses $O(L)$ seed bits since $t(a+1)=O(L)$.

Fix the choices of all preceding stages and consider one retained set
of size $m>2^{\lceil a/2\rceil}$. Choose an integer $0\le d\le a$ such
that $\mu=m2^{-d}$ lies in $[2^{a/4},2^{1+a/4}]$. Let $I_v$ indicate
that the first $d$ bits of vertex $v$ are zero, and let $X$ be their sum
over the retained set. Under fully independent bits these indicators
are independent with mean $p=2^{-d}$. In the expansion of
$\mathbb E[(X-\mu)^t]$, a tuple in which some vertex occurs exactly
once contributes zero. A remaining tuple has $j\le t/2$ distinct
vertices and contributes at most $(2p)^j$ in absolute value, since
$\mathbb E|I_v-p|^h\le 2p$ for all integers $h\ge 1$. There are at
most $m^j t^t$ tuples with $j$ distinct vertices. Hence the moment
under fully independent bits is at most
$2t^t(2\mu)^{t/2}$.

The Fourier coefficients of $I_v$ on the binary cube have absolute
sum one, by expanding the product of the $d$ indicators of a zero bit.
Thus the coefficients of $I_v-p$ have absolute sum at most two.
The coefficients of each product of $t$ centered indicators have
absolute sum at most $2^t$, including when vertices repeat. Under an
$\varepsilon$-biased distribution, its expectation therefore differs
from that under independence by at most $2^t\varepsilon$. Summing over
the $m^t$ tuples and applying Markov's inequality for even $t$ gives
\[
 \Pr[|X-\mu|\ge\mu]
 \le 2(2t^2/\mu)^{t/2}
      +\varepsilon(2m/\mu)^t
 \le 2^{-5L}
\]
for all sufficiently large $N$. Indeed, $a>256\ell$ and
$\log_2 t\le \ell+7$ give
$\log_2(\mu/(2t^2))\ge a/6$, while $t\ge128L/a$.
On the complementary event the prefix class contains at least one and
fewer than $2\mu\le 2^{\lceil a/2\rceil}$ vertices. The least occupied
color class is contained in this prefix class. A set already of size at
most $2^{\lceil a/2\rceil}$ stays within the bound. A union bound gives
a seed that works for all the retained sets at once. Since this holds
for each fixed successful prefix of seeds, a successful sequence exists.

There are $O(\ell)$ stages. The final bound $K=2^{256\ell}$ is
polynomial in $L$. The full seed has $O(L\ell)$ bits. Encode the
tuple of colors as a mixed-radix integer $c(v)\in\{0,\ldots,C-1\}$ in
lexicographic order. Each radix is at most $L+1$, so
$\log_2 C=O(\ell^2)$.

Apply the coloring to the nonempty feasible midpoint sets of the graph.
Put $B=4C+4$ and, at a positive layer $i$, set
\[
 V(v)=(1+c(v))B^{\nu_2(i)}.
\]
Take $V=0$ in layer zero. Fix an endpoint pair with pivot of valuation
$h$. For each $j\ge 1$ there are at most $2^j$ internal layers of
valuation $h-j$. Consequently their total contribution, summed over
all lower valuations along any path, is at most
\[
 C B^h\sum_{j\ge 1}(2/B)^j
 =\frac{2C}{B-2}B^h < B^h.
\]
An increase of one in pivot color costs $B^h$. Any minimum-$V$ path
must therefore use a feasible midpoint of the least occupied color.
There are at most $K$ such vertices. Finally $h\le L$, and a path has
at most $N$ vertices, giving bit length $O(L\ell^2)$. The colors and
the integer power can be evaluated within this space by regenerating
the relevant bits and using binary arithmetic.
\end{proof}
\section{Frontiers and the lists of endpoint pairs}

Fix successful initial weights $V$ from Lemma~\ref{lem:raw}. A \emph{window}
$I=(x,y)$ is a reachable ordered pair with at least one internal layer.
For $D\ge1$, give layer $t>0$ rank
\[
                   r_D(t)=\lfloor \nu_2(t)/D\rfloor .
\]
The \emph{frontier} $\fr_D(I)$ is the nonempty set of internal layer indices
of largest rank $r_D$ in the window. The $D$-skeleton of a path in $I$
lists the vertices it uses on these layers, in order. For a vector $W$
of weights, let $\calF_D(I;W)$ be the set of $D$-skeletons of $W$-optimal paths.
Write $f_D(I;W)=|\calF_D(I;W)|$. In particular, $f_1(I;V)\le K$.

Put $D_i=2^i$, for $0\le i\le s$, where $s$ is largest with
$D_s\le L/\ell^2$. Thus $s=O(\ell)$ and $D_s=\Theta(L/\ell^2)$.
We pass between successive frontiers as follows. Put $D=D_{i-1}$.
The binary valuations in a $2D$-frontier belong to one
band of $2D$ consecutive valuations. This band splits into a lower
half and an upper half of $D$ valuations each.
\begin{itemize}
\item First expose the $D$-skeleton of the entire window $I$.
If the greatest internal valuation belongs to the lower half, this
already specifies the $2D$-skeleton.
\item Otherwise the first exposure gives all upper-half vertices.
Between consecutive specified vertices, including the original endpoints,
expose the $D$-skeleton of each gap containing a lower-half frontier
layer. These exposures give precisely the remaining vertices of the
$2D$-skeleton.
\end{itemize}
Call each window on which this procedure asks for a $D$-skeleton a
\emph{site}. The site endpoints in the second step depend on the first
answer. On a fixed path the sites' $D$-frontiers partition $\fr_{2D}(I)$.
Site intervals can be nested, but their frontier layer sets are disjoint.

For each $i\in[s]$ define
\[
 a_i=D_i\lceil\log_2(K+1)\rceil+\ell,\qquad
 b_i=\lceil C_*\ell a_i\rceil,\qquad m_i=\lceil 16 L/b_i\rceil ,
\]
where $C_*$ is a sufficiently large absolute constant fixed below.
The integer $b_i$ controls the range of a perturbation at step $i$.
The number $m_i$ is how many ambiguous sites a witness must exhibit.

Define finite catalogues $\calC_s,\ldots,\calC_0$ before choosing any
perturbations. A member of $\calC_i$ is a list of windows with pairwise
disjoint $D_i$-frontiers. Start $\calC_s$ with the singleton lists of all
windows. Given $\calC_i$, take a list $J\in\calC_i$. Form the union of
the possible sites obtained from its windows using any $V$-minimum path
in each window. Put in $\calC_{i-1}$ all lists of length at most $m_i$
drawn from this union, in any order, whose $D_{i-1}$-frontiers are
pairwise disjoint. Include the empty list. Allowing separately feasible
sites from incompatible paths only enlarges the catalogue.

\begin{lemma}\label{lem:catalogues}
For $C_*$ sufficiently large and $N$ sufficiently large, every catalogue
has size at most $2^{4L}$.
\end{lemma}
\begin{proof}
Fix a parent list and consider how to specify one possible site. Its
parent window can be named in $O(\ell)$ bits, since every list has
length $O(L)$. In a band of $D_i$ binary valuations the frontier contains
at most $2^{D_i}-1$ positions. Indeed the positions are multiples of
the band's lowest power of two between consecutive multiples of the
next higher band's power. A site can therefore name its endpoint layers
in $O(D_i+1)$ bits, allowing the two original endpoints as well.

To name the vertex in one of those layers, descend the binary pivot
recursion from the two known endpoints of the parent. At each pivot
there are at most $K$ choices under $V$. Once the vertex has been
chosen, the target layer determines which side to enter. The valuations
strictly decrease and the target belongs to the same band as the first
pivot, so at most $D_i$ vertex choices are needed. The same argument
names the other endpoint. This is a counting argument, so the names may
be ranks in lists defined by $V$-optimality. A total of $c a_i$ bits
per site suffices for an absolute $c$, including the length and ordering
overhead for specifying a list of up to $m_i$ sites by padding.

It follows that
\[
 \log_2|\calC_j|\le 2L+c\sum_{i=j+1}^{s}m_i a_i
 \le 2L+\frac{16cLs}{C_*\ell}
             +c\sum_{i=1}^{s}a_i .
\]
The last sum is $O(D_s\log(K+1)+s\ell)=O(L/\ell+\ell^2)=o(L)$.
Choose $C_*$ large enough that the middle term is at most $L$.
\end{proof}

\section{Controlling ambiguity on all the listed windows}

We construct a new additive objective $U_i$ at each step. Write
$W_0=(V)$ and $W_i=(V,U_1,\ldots,U_i)$ with lexicographic minimization.
We will inductively establish
\begin{equation}\label{eq:inv}
     \prod_{I\in J} f_{D_i}(I;W_i)\le 2^L
                  \qquad\text{for every }J\in\calC_i .
\end{equation}
For $i=0$ this follows from $f_1(I;V)\le K$ and
$m_1\log_2 K\le L$ for sufficiently large $N$.

\subsection{Weights at one step}

Fix $i\ge1$, write $D=D_{i-1}$ and $q=2^D$, and choose the least prime
$p>2^{b_i}$. Take a power of two $A$ between $16p q^2$ and $32p q^2$.
The rank $r_D(t)$ is also the exponent of $q$ dividing $t$.
There are at most $H=1+\lfloor L/D\rfloor$ ranks.
Partition them in increasing order into consecutive batches of at most
$\Delta$ ranks. The choice of $\Delta$ will be made in the space calculation.
Independently in each batch $g$, choose $h_g\in\F_p^N$ using
Lemma~\ref{lem:bias} with $\varepsilon=2^{-32L}$. Give a vertex $v$
whose rank is in batch $g$, at position $j$ counting from zero, the weight
\[
                    U_i(v)=A^j[h_g(v)]_p,
\]
where $[\cdot]_p$ is the representative in $\{0,\ldots,p-1\}$.
All objectives assign zero to layer zero.

\begin{lemma}[Simultaneous ties]\label{lem:ties}
Condition on $W_{i-1}$. Fix $m$ sites whose $D$-frontiers are disjoint.
At each site $I$ fix two distinct skeletons from $\calF_D(I;W_{i-1})$.
The probability that at every site both fixed skeletons occur on
$W_i$-optimal paths is at most $p^{-m}+2H\varepsilon$.
\end{lemma}
\begin{proof}
Consider one site with greatest internal rank $h$ at position $j$ in
batch $g$. Condition on the seeds of earlier batches. For a fixed
skeleton $S$, let $c_S$ be the minimum earlier-batch $U_i$-cost over
the $W_{i-1}$-optimal paths in $I$ with skeleton $S$.
It is fixed before $h_g$ is chosen. There are no internal layers of
later batches. The least total $U_i$-cost in this class equals
\[
      c_S + A^j \sum_{v\in S}[h_g(v)]_p+\eta_S,
                         \qquad 0\le\eta_S<A^j/8 .
\]
For the error bound, there are fewer than $q^{t+1}$ internal layers
of rank $h-t$ for $1\le t\le h$, since the interval contains no internal
multiple of $q^{h+1}$. Thus the portion from smaller ranks of the
current batch is bounded, on every path, by
\[
    A^j p\sum_{t\ge1} q^{t+1} A^{-t}
      = A^j \frac{p q^2}{A-q}<A^j/8 .
\]
Minimization over paths preserves this bound.

If $S$ and $T$ both occur on $W_i$-optimal paths, their class minima
are equal. The integer $\sum_{v\in S}[h_g(v)]_p-\sum_{v\in T}[h_g(v)]_p$
lies within $1/8$ of $(c_T-c_S)/A^j$. There is at most one such integer,
fixed under the conditioning. Reducing it modulo $p$ gives an affine
constraint, or the event is impossible. Its coefficient row is nonzero
because different skeletons differ at some layer. Within this batch,
distinct sites use disjoint frontier layers, hence their coefficient
rows have disjoint supports and are linearly independent.

If $m_g$ sites belong to batch $g$, the conditional probability of the
required constraints is at most $p^{-m_g}+\varepsilon$. Whether the ties
occur in one batch depends only on that batch and its predecessors.
Iterating conditional probabilities gives an upper bound
$\prod_g(p^{-m_g}+\varepsilon)$. Its difference from $p^{-m}$ is at most
$(1+\varepsilon)^H-1\le2H\varepsilon$ for sufficiently large $N$.
\end{proof}

\subsection{The joint induction}

\begin{proposition}\label{prop:ind}
If~\eqref{eq:inv} holds at step $i-1$, the batch seeds for $U_i$ can
be chosen so that it holds at step $i$.
\end{proposition}
\begin{proof}
For a list $J$ of exactly $m_i$ sites in $\calC_{i-1}$, the number of ways
to specify an ordered pair of old-optimal skeletons at every site is at most
\[
                   \left(\prod_{I\in J}f_D(I;W_{i-1})\right)^2\le 2^{2L},
                   \qquad D=D_{i-1}.
\]
This use of the product bound is essential. It costs only $2L$ bits to
name the two entire tuples, even though the list can have many sites.
Consider choices that differ at every site. By Lemmas~\ref{lem:catalogues}
and~\ref{lem:ties}, the probability that any such choice survives
everywhere is at most
\[
          2^{4L}\,2^{2L}\bigl(p^{-m_i}+2H\varepsilon\bigr)<1,
\]
using $m_i b_i\ge16L$. Fix seeds avoiding this event.

Fix a parent list $J\in\calC_i$. Form a decision tree whose leaves are
the tuples of $D_i$-skeletons obtained by choosing a $W_i$-optimal path
in each listed window. Different windows may choose paths independently.
Use list order and the two-step frontier decomposition above, processing
gaps in layer order. At each site branch on the different $D$-skeletons
that admit a completion consistent with the prefix. This prescription
gives each full tuple one leaf, since the sequence of site questions
is determined by the earlier answers. Along a given traversal the
site frontiers are disjoint. This follows from the partition inside
each window and the disjoint parent frontiers.

If a traversal had $m_i$ sites with more than one available answer,
take exactly $m_i$ of them, with the endpoints they have on this traversal.
They form a list in $\calC_{i-1}$. Indeed, every $W_i$-optimal path
is $V$-minimum, and its site placements therefore occur among those used
to construct the catalogue. At each branching site at least two skeletons
occur on $W_i$-optimal paths of that site, by optimal subpath replacement.
The alternative at an earlier site need not preserve later endpoints.
The excluded event only asks that the two alternatives exist separately
at each of the recorded sites. This contradiction shows that every
traversal has fewer than $m_i$ branching sites.

Choose an answer uniformly from the available answers at each decision
node. For any fixed leaf, the list of all its branching sites also belongs
to $\calC_{i-1}$. Its probability is the reciprocal of the product of
its branching factors. Those factors are bounded by the respective
$f_D(I;W_{i-1})$, so by induction every leaf has probability at least
$2^{-L}$. The tree has at most $2^L$ leaves. Its leaf count is precisely
the product in~\eqref{eq:inv} for the parent list. This proves the
inductive claim for every parent simultaneously.
\end{proof}

\section{Isolation and space}

At the final scale the catalogue includes every singleton, so
$f_{D_s}(I;W_s)\le2^L$ for every window. Choose a prime $P>N^{10}$
of $O(L)$ bits and one field element $z\in\F_P$. For distinct old-optimal
skeletons $S,T$ at any one window, the polynomial
\[
           \sum_{v\in S}z^{\id(v)}-\sum_{v\in T}z^{\id(v)}
\]
is nonzero and has degree at most $N$. The number of comparisons is at
most $N^2 2^{2L}$, so some $z$ makes all these differences nonzero.
For the last objective give $v$ in positive layer $t$ weight
\[
                   B^{r_{D_s}(t)}[z^{\id(v)}]_P ,
                    \qquad B=16NP .
\]
On a window of maximum rank $h$, a unit gap between frontier sums
dominates all contributions of lower ranks, which total at most
$NP B^{h-1}$ when $h>0$. The winning skeleton is unique, and its
gap intervals have smaller maximum rank. By induction their optimal paths
are unique. A pair with no internal layer has at most one path by
simplicity. Thus the complete lexicographic objective is min-isolating.

\paragraph{Size of the weights and seed.}
The initial objective costs $O(L\ell^2)$ bits by Lemma~\ref{lem:raw}.
At step $i$, use $u_i=b_i+2D_{i-1}+7$ as an integer upper bound on
$\log_2 A$, using Bertrand's postulate. We have
\[
 u_i=O(D_{i-1}\ell^2),\qquad b_i=O(L),\qquad
 H=O(L/D_{i-1}).
\]
Each batch seed costs $O(L)$ by Lemma~\ref{lem:bias}. For $\Delta$ ranks
per batch, all seeds at this step and the bit length of its path cost
therefore total
\[
             O\bigl(L\lceil H/\Delta\rceil+L+\Delta u_i\bigr).
\]
Take $\Delta$ to be $\lceil\sqrt{HL/u_i}\rceil$ clipped to $[1,H]$.
The displayed quantity is $O(L+u_i+\sqrt{H L u_i})=O(L\ell)$.
All $O(\ell)$ steps together cost $O(L\ell^2)$. The last objective
has $O(L/D_s)=O(\ell^2)$ radix positions of $O(L)$ bits. All the parameters
admit integer computation using ceilings or integer powers of two.

Encode the objectives in mixed radix, taking each radix greater than
an explicit upper bound on the maximum path cost of its coordinate.
Multiplying the maximum possible vertex weight for that coordinate by
$N$ suffices. Summing the encoded vertex weights then implements
lexicographic path comparison without a carry. The resulting bit length
is bounded by the sum already computed. Trial division for primes and
field polynomials, evaluation of the small-bias coordinates, and
ordinary integer arithmetic fit the same space.

\paragraph{Turning isolation into an unambiguous decision.}
Enumerate all seed tuples in fixed order until the first min-isolating
tuple passes. The weighted reachability checker of van Melkebeek and
Prakriya~\cite[Lemma 3.5]{vMP} uses space $O(\log N+\log(W+1))$ for an
upper bound $W$ on the path weights. For completeness, in the layered
case one can implement the checker as follows. Count all entered vertices
when weighing paths. At layer distance $j$ store the number of reachable
pairs and the sum of their minimum weights. For $j=0$ these numbers are
$N$ and zero. Suppose the census is correct and all the minima at distance
$j$ are unique. To look up any of these distances, scan all pairs at that
layer distance, guessing either to skip a pair or exhibit a path for it.
Abort unless the stored count and sum match. The count forces every
reachable pair to be represented. The sum then forces every exhibited path
to be minimum. Using a unique encoding of path guesses leaves exactly one
returning branch, which also knows the requested distance or its infinity.

For each distance-$(j+1)$ pair, scan the last edge and use these lookups
to find the best predecessor. Two predecessors tied for a finite minimum give a unique
return reporting a failed assignment. Otherwise accumulate the next
census and continue. This uses a constant number of counters of
$O(L+\log(W+1))$ bits. One can also record the designated pair's reachability,
handling same-layer and backward queries directly. Thus even the search
through failed assignments has a unique continuing branch. The enumeration
terminates by Proposition~\ref{prop:ind} and the final isolation argument.
This proves Theorem~\ref{thm:main}.

\begin{thebibliography}{9}
\bibitem{vMP} Dieter van Melkebeek and Gautam Prakriya.
\emph{Derandomizing isolation in space-bounded settings}.
SIAM Journal on Computing 48(3), 979--1021, 2019.
\href{https://doi.org/10.1137/17M1130538}{doi:10.1137/17M1130538}.
\bibitem{Hoza} William M. Hoza.
\emph{Deterministic, oblivious isolation for space-bounded computation requires large weights}.
ECCC TR26-124, 2026. \url{https://eccc.weizmann.ac.il/report/2026/124/}
\bibitem{RA} Klaus Reinhardt and Eric Allender.
\emph{Making nondeterminism unambiguous}.
SIAM Journal on Computing 29(4), 1118--1131, 2000.
\href{https://doi.org/10.1137/S0097539798339041}{doi:10.1137/S0097539798339041}.
\end{thebibliography}
\end{document}
